% =====================================================================
%  MDL-Optimal Schema Selection for Bit-Uniform Relational Storage
%  Formal technical note -- self-contained LaTeX source
% =====================================================================
\documentclass[10pt]{article}

\usepackage{amsmath}
\usepackage{amssymb}
\usepackage{amsthm}
\usepackage{booktabs}
\usepackage{algorithm}
\usepackage{algpseudocode}
\usepackage{listings}
\usepackage{xcolor}
\usepackage{hyperref}
\usepackage{cleveref}
\usepackage{geometry}
\usepackage{natbib}

\geometry{margin=1in,letterpaper}

\hypersetup{
  colorlinks=true,
  linkcolor=blue!45!black,
  citecolor=blue!45!black,
  urlcolor=blue!45!black
}

% ---- Theorem environments ------------------------------------------
\newtheorem{theorem}{Theorem}
\newtheorem{lemma}[theorem]{Lemma}
\newtheorem{corollary}[theorem]{Corollary}
\newtheorem{proposition}[theorem]{Proposition}
\theoremstyle{definition}
\newtheorem{definition}{Definition}
\theoremstyle{remark}
\newtheorem{remark}{Remark}

% ---- Listings ------------------------------------------------------
\lstset{
  basicstyle=\ttfamily\footnotesize,
  keywordstyle=\color{blue!70!black}\bfseries,
  commentstyle=\color{green!45!black},
  stringstyle=\color{red!55!black},
  frame=single,
  breaklines=true,
  showstringspaces=false,
  columns=fullflexible
}

% ---- Shorthands ----------------------------------------------------
\newcommand{\Vcal}{\mathcal{V}}
\newcommand{\Ccal}{\mathcal{C}}
\newcommand{\Tcal}{\mathcal{T}}
\newcommand{\Lmodel}{L_{\mathrm{model}}}
\newcommand{\Ldata}{L_{\mathrm{data}}}
\newcommand{\bits}{\{0,1\}^*}
\newcommand{\eqL}{\equiv_{\mathcal{L}}}
\DeclareMathOperator*{\argmin}{arg\,min}

\title{MDL-Optimal Schema Selection for Bit-Uniform Relational Storage}
\author{Anonymous}
\date{}

\begin{document}

\maketitle

% =====================================================================
\begin{abstract}
Modern columnar storage engines increasingly defer type interpretation to
read time, storing raw $64$-bit NaN-boxed words in a bit-uniform layout.
This deferral enables schema evolution and polymorphic columns, but it
raises a fundamental question: which type interpretation should a reader
apply to an opaque column of bits? We formalize this schema-selection
problem as a minimum description length (MDL) optimization. Given a
column $\Ccal$ of $n$ raw $64$-bit words and a finite candidate type set
$\Tcal$, we seek the encoding $\tau \in \Tcal$ minimizing the two-part
code $L(\Ccal,\tau) = \Lmodel(\tau) + \Ldata(\Ccal \mid \tau)$, where
$\Lmodel$ charges for the type tag and metadata and $\Ldata$ charges for
the per-value code under $\tau$. We prove that the MDL-optimal encoding is
unique up to information-lattice equivalence, that niche-filling incurs no
description-length overhead, and that the selection algorithm runs in
$O(|\Tcal| \cdot n)$ time. A corollary establishes that read-time
reinterpretation is safe if and only if the candidate encodings are
information-lattice equivalent.
\end{abstract}

% =====================================================================
\section{Introduction}\label{sec:intro}

Relational storage engines have, for decades, committed to a fixed schema
at write time: each column is assigned a concrete physical type, and every
value stored in that column is encoded according to that type. This
discipline simplifies the query executor but rigidifies the on-disk
representation. When the schema evolves---an integer column must now hold a
floating-point sentinel, or a string column acquires nulls---storage
engines resort to costly migrations, rewrite-on-alter, or auxiliary shadow
columns; variant indexes~\citep{oneil1997} anticipated this flexibility at the index level decades ago. A growing line of systems instead stores \emph{bit-uniform}
columns of raw $64$-bit words and defers type interpretation to read time,
recovering the concrete type from the bit pattern itself through
\emph{NaN-boxing}~\citep{wingo2011}, a representation pioneered by
JavaScript runtimes and now ubiquitous in dynamic-language
implementations.

Schema-on-read demands a principled objective. Without one, the choice of
type interpretation degenerates into a catalog of ad hoc heuristics: ``use
\texttt{i32} when all values fit, otherwise fall back to \texttt{f64}.''
Such heuristics are fragile, mutually inconsistent, and invisible to the
query planner. The minimum description length (MDL)
principle~\citep{rissanen1978,grunwald2004,barron1998} supplies a unified
objective: among a family of candidate encodings, choose the one that
minimizes the total number of bits required to transmit both the model
(the type and its metadata) and the data (the encoded values). MDL is
model-class-agnostic, statistically consistent, and---crucially for
us---computable in a single linear pass over the column.

This note makes the following contributions. First, we give a formal
specification of the schema-selection problem for bit-uniform storage
(\cref{sec:problem}) and a closed-form MDL objective whose model and data
parts decompose cleanly across the NaN-boxed word layout
(\cref{sec:objective}). Second, we present \textsc{Schema-Select}, a
linear-time-per-candidate algorithm with a tight $O(|\Tcal| \cdot n)$
complexity bound (\cref{sec:algorithm}). Third, we connect the selection
problem to Verd\'u's information lattice~\citep{verdu} and prove that
MDL-optimal encodings are unique up to information-lattice equivalence
(\cref{sec:lattice}); the accompanying corollary gives the precise
condition under which read-time reinterpretation is semantically safe.
Fourth, we treat variable-length type tags via Huffman
coding~\citep{huffman1952} and a fully worked numerical example
(\cref{sec:tags,sec:worked}).

The remainder of the note proceeds as follows. \Cref{sec:prelim} reviews
NaN-boxing, niche-filling, and the MDL principle. \Cref{sec:problem}
states the problem. \Cref{sec:objective,sec:algorithm} develop the
objective and the algorithm. \Cref{sec:lattice} establishes the
information-lattice equivalence theorem.
\Cref{sec:tags,sec:worked,sec:impl} address practical concerns, and
\Cref{sec:concl} concludes. The development is self-contained; we assume
only familiarity with elementary information theory~\citep{coverthomas}
and the relational model.

% =====================================================================
\section{Preliminaries}\label{sec:prelim}

\subsection{Notation}\label{sec:notation}

Let $\Vcal$ denote the universe of values a relational attribute may
assume---integers, IEEE-754 doubles, UTF-8 strings, nulls, and tagged
variants thereof. A \emph{column} of length $n$ is a finite sequence
$\Ccal = (v_1, v_2, \ldots, v_n)$ with $v_i \in \Vcal$. A \emph{type
interpretation} $\tau \in \Tcal$ is a pair $(e_\tau, d_\tau)$ where
$e_\tau : \Vcal \to \bits$ is an injective \emph{encoder} and
$d_\tau : \bits \rightharpoonup \Vcal$ is a partial \emph{decoder} with
$d_\tau \circ e_\tau = \mathrm{id}$ on $\Vcal$. The candidate type set
$\Tcal$ is finite and fixed by the storage engine; throughout this note
we take $\Tcal = \{\textsc{f64}, \textsc{i32}, \textsc{\&str},
\textsc{nullable}, \textsc{variant}\}$ as a running example. We write
$\bits$ for the set of finite binary strings and $\{0,1\}^{64}$ for the
set of $64$-bit words. The length of a string $x$ is $|x|$.

\subsection{IEEE-754 NaN-boxing recap}\label{sec:nanbox}

An IEEE-754 binary64 double~\citep{coverthomas} is a $64$-bit word
partitioned into a sign bit, an $11$-bit exponent, and a $52$-bit
mantissa. A value is a \emph{NaN} precisely when all eleven exponent bits
are set and the mantissa is nonzero, yielding $2^{53} - 2$ distinct NaN
bit patterns (counting the sign). Of these, the canonical single NaN
specified by the standard is but one; the remaining patterns form a vast
\emph{payload space} that the storage engine may exploit. The
\emph{NaN-boxing} technique~\citep{wingo2011} co-opts this payload space
as a type-tagged container: the upper $16$ bits (sign and exponent)
encode a type tag, while the lower $48$ bits carry a type-specific
payload. \Cref{tab:nanbox} gives the partition used throughout this note.

\begin{table}[h]
\centering
\caption{NaN-boxed $64$-bit word partition. The upper $16$ bits serve as
a type tag; the lower $48$ bits carry the payload. Finite doubles occupy
the non-NaN region and require no tag.}\label{tab:nanbox}
\begin{tabular}{@{}llll@{}}
\toprule
Upper $16$ bits (hex) & Class & Type tag $\tau$ & Payload ($48$ bits) \\
\midrule
\texttt{0x0000}--\texttt{0x7FEF} & finite \textsc{f64}    & \textsc{f64}     & IEEE-754 double value (full $64$ bits) \\
\texttt{0x7FF0}                  & $+\infty$               & \textsc{f64}     & --- \\
\texttt{0x7FF8}                  & quiet NaN               & \textsc{i32}     & $32$-bit signed integer (low $32$, zero-extended) \\
\texttt{0x7FF9}                  & quiet NaN               & \textsc{\&str}   & $48$-bit offset into string heap \\
\texttt{0x7FFA}                  & quiet NaN               & \textsc{null}    & --- (empty payload) \\
\texttt{0x7FFB}                  & quiet NaN               & \textsc{variant} & $8$-bit variant tag $+\;40$-bit payload \\
\texttt{0xFFF0}                  & $-\infty$               & \textsc{f64}     & --- \\
\texttt{0xFFF8}--\texttt{0xFFFF} & sign-bit NaN            & reserved         & future extensions \\
\bottomrule
\end{tabular}
\end{table}

The essential observation is that finite doubles---the overwhelming
majority of values in numerical columns---pay no tagging overhead at all:
they live in the non-NaN region and are stored verbatim. Only the
``exotic'' values (integers, strings, nulls, variants) occupy the NaN
region, where they are tagged at the cost of the upper $16$ bits. This
asymmetry is what makes NaN-boxing a compelling substrate for MDL
analysis: the description length of a column is dominated by the data
part, and the data part is, for homogeneous numerical columns, exactly
$n \cdot 64$ bits.

\subsection{Niche-filling}\label{sec:niche}

\emph{Niche-filling}~\citep{rustlayout,swiftabi} is the systematic
allocation of unused bit patterns within a representation to alternative
types. Rust's enum layout and Swift's resilient ABI both exploit the
observation that pointer values on $64$-bit platforms never use the upper
$16$ bits of a $64$-bit word (the user-space address space is $48$ bits),
freeing those bits to discriminate enum variants. The NaN payload is a
special case: the niche is the set of NaN bit patterns, and the
alternative types are the non-double value classes. Niche-filling is
attractive precisely because it is \emph{free}: the bits it occupies
would otherwise be wasted. We make this intuition precise in
\cref{lem:niche}: niche-filling preserves description length, in the
sense that storing a tagged value in the NaN payload adds no model cost
beyond the type tag itself. This is the formal basis for treating the
NaN-boxed layout as a bit-uniform substrate over which MDL selection
operates.

\subsection{The MDL principle}\label{sec:mdl}

The minimum description length principle~\citep{rissanen1978,grunwald2004}
holds that, given a family of candidate models $\Tcal$ and observed data
$\Ccal$, one should select the model $\tau^\star \in \Tcal$ minimizing the
\emph{two-part code} length
\begin{equation}\label{eq:twopart}
L(\Ccal, \tau) \;=\; \Lmodel(\tau) \;+\; \Ldata(\Ccal \mid \tau),
\end{equation}
where $\Lmodel(\tau)$ is the cost, in bits, of describing the model (here,
the type tag and its metadata) and $\Ldata(\Ccal \mid \tau)$ is the cost
of describing the data given the model. Rissanen's original
formulation~\citep{rissanen1978} charged $\Lmodel$ with the logarithm of
the model index; the \emph{refined MDL} of \citet{barron1998} replaces
this with the normalized maximum likelihood (NML) code, which accounts for
the complexity of the model class itself. Both formulations trace
ultimately to Kolmogorov complexity~\citep{livitanyi}, which charges a
program for the length of its shortest description. For our purposes the
distinction is benign: the candidate type set $\Tcal$ is fixed and finite,
so the model cost reduces to a constant per-type overhead. We adopt the
two-part code of \cref{eq:twopart} throughout, noting that the
information-lattice uniqueness result of \cref{sec:lattice} holds
verbatim under NML because NML refines, rather than replaces, the
two-part ranking within a fixed model class.

% =====================================================================
\section{Problem statement}\label{sec:problem}

We now state the schema-selection problem formally. Let $\Ccal =
(v_1, \ldots, v_n)$ be a column of raw $64$-bit NaN-boxed words supplied
by the storage engine at load time, and let $\Tcal$ be the finite set of
candidate type interpretations admitted by the engine. The decision
variable is a per-column type assignment $\tau \in \Tcal$. The objective
is to minimize the total description length $L(\Ccal, \tau)$ of
\cref{eq:twopart} over $\Tcal$:
\begin{equation}\label{eq:objective}
\tau^\star \;=\; \argmin_{\tau \in \Tcal} \; L(\Ccal, \tau) \;=\; \argmin_{\tau \in \Tcal} \;\bigl[\,\Lmodel(\tau) + \Ldata(\Ccal \mid \tau)\,\bigr].
\end{equation}
The assignment is per-column and per-load: different columns in the same
relation may receive different types, and the same column may be
reinterpreted across loads if the data drifts. We restrict attention to
\emph{lossless} encodings, meaning $d_\tau \circ e_\tau = \mathrm{id}$;
lossy compression is orthogonal and composes with our scheme. The problem
is well-posed because $\Tcal$ is finite and $L$ is real-valued, so a
minimizer exists. The remainder of the note characterizes the minimizer
(\cref{sec:objective,sec:lattice}), gives an efficient algorithm for
computing it (\cref{sec:algorithm}), and demonstrates its behavior on
concrete data (\cref{sec:worked}).

% =====================================================================
\section{The MDL objective}\label{sec:objective}

We instantiate the abstract two-part code of \cref{eq:twopart} for the
NaN-boxed setting. The model cost decomposes into three additive
terms: a \emph{type tag} identifying $\tau$ within $\Tcal$, an
\emph{index structure} cost for any auxiliary arrays the encoding
requires, and a \emph{metadata} cost for per-type parameters such as
string-heap length or variant tag table. Concretely,
\begin{equation}\label{eq:model}
\Lmodel(\tau) \;=\; \underbrace{\lceil \log_2 |\Tcal| \rceil}_{\text{type tag}} \;+\; \underbrace{I(\tau)}_{\text{index}} \;+\; \underbrace{M(\tau)}_{\text{metadata}},
\end{equation}
where $I(\tau)$ is the size, in bits, of any offset or indirection array
(e.g., a $48$-bit offset table for \textsc{\&str}, over which string
search proceeds by the methods of \citet{baeza1992}) and $M(\tau)$ is the
size of per-type parameters (e.g., the variant tag enumeration for
\textsc{variant}). For the homogeneous numeric types \textsc{f64} and
\textsc{i32} both $I$ and $M$ vanish, leaving only the tag. The data cost
is a per-value sum,
\begin{equation}\label{eq:data}
\Ldata(\Ccal \mid \tau) \;=\; \sum_{i=1}^{n} \ell_\tau(v_i),
\qquad \ell_\tau(v) \;=\; |e_\tau(v)|,
\end{equation}
where $\ell_\tau(v)$ is the length of the encoding of $v$ under $\tau$.
For a finite double under \textsc{f64}, $\ell_{\textsc{f64}}(v) = 64$. For
an integer under \textsc{i32}, $\ell_{\textsc{i32}}(v) = 32$ (the low
$32$ bits of the NaN-boxed word, the upper $32$ being a fixed tag). For a
value outside the domain $\mathrm{dom}(\tau) \subsetneq \Vcal$, we adopt
the \emph{escape convention}: the value is stored as an opaque $64$-bit
word, so $\ell_\tau(v) = 64$ whenever $v \notin \mathrm{dom}(\tau)$. This
convention keeps every candidate $\tau$ admissible for every column,
routing the cost of heterogeneity into $\Ldata$ rather than into an
artificial validity constraint.

\begin{definition}[Bit-uniform encoding]\label{def:bituniform}
A type interpretation $\tau$ is \emph{bit-uniform} if every value
$v \in \mathrm{dom}(\tau)$ is encoded into a $64$-bit word, i.e.,
$e_\tau : \mathrm{dom}(\tau) \to \{0,1\}^{64}$, and the storage engine
allocates exactly $64$ bits per value on disk regardless of $\tau$.
\end{definition}

Bit-uniformity is the physical invariant that makes MDL selection
non-trivial: because every value occupies the same number of physical
bits, the only lever for reducing description length is the choice of
$\tau$, which controls how many of those bits are \emph{semantic}
(payload) versus \emph{structural} (tag and padding). A non-bit-uniform
engine would already have committed to a variable-length layout, and the
selection problem would collapse to conventional compression analysis.

\begin{definition}[Description length]\label{def:dl}
Given a column $\Ccal = (v_1, \ldots, v_n)$ and a bit-uniform type
interpretation $\tau$, the \emph{description length} of $\Ccal$ under
$\tau$ is
\[
L(\Ccal, \tau) \;=\; \Lmodel(\tau) + \Ldata(\Ccal \mid \tau)
\;=\; \lceil \log_2 |\Tcal| \rceil + I(\tau) + M(\tau) + \sum_{i=1}^{n} \ell_\tau(v_i).
\]
\end{definition}

As a concrete illustration, consider a column of $n = 10^6$ values known
to be $32$-bit integers. Under \textsc{f64} the data cost is
$10^6 \cdot 64 = 6.4 \times 10^7$ bits; under \textsc{i32} it is
$10^6 \cdot 32 = 3.2 \times 10^7$ bits. The model cost is identical
($\lceil \log_2 5 \rceil = 3$ bits) in both cases, so the MDL objective
prefers \textsc{i32} by a margin of $3.2 \times 10^7$ bits---a factor of
two in storage. Under \textsc{\&str}, by contrast, every integer is an
out-of-domain escape costing $64$ bits plus the offset-table index cost
$I(\textsc{\&str}) \approx 48 \cdot 10^6$ bits, yielding a total roughly
double the \textsc{f64} cost. The objective thus recovers the
``obvious'' choice by construction, without an explicit rule.

% =====================================================================
\section{The selection algorithm}\label{sec:algorithm}

\Cref{alg:schema-select} presents \textsc{Schema-Select}, a direct
instantiation of the argmin in \cref{eq:objective}. The algorithm scans
the column once per candidate type, accumulating the per-value cost
$\ell_\tau(v_i)$, and retains the candidate of minimum total length. The
inner loop admits a single-pass streaming implementation: each word is
classified by its upper $16$ bits (a $O(1)$ table lookup into the NaN-box
partition of \cref{tab:nanbox}), and $\ell_\tau(v_i)$ is read from a
precomputed cost table indexed by $(\tau, \text{class})$. No sorting,
hashing, or auxiliary data structures are required.

\begin{algorithm}[h]
\caption{\textsc{Schema-Select}}\label{alg:schema-select}
\begin{algorithmic}[1]
\Require Column $\Ccal = (v_1, \ldots, v_n)$ as raw $64$-bit words; candidate types $\Tcal$
\Ensure MDL-optimal type interpretation $\tau^\star$
\State $L^\star \gets +\infty$;\quad $\tau^\star \gets \bot$
\For{$\tau \in \Tcal$}
    \State $L_m \gets \textsc{ModelCost}(\tau)$ \Comment{$\lceil \log_2 |\Tcal| \rceil + I(\tau) + M(\tau)$}
    \State $L_d \gets 0$
    \For{$i = 1$ \textbf{to} $n$}
        \State $c \gets \textsc{Classify}(v_i)$ \Comment{upper $16$-bit lookup}
        \State $L_d \gets L_d + \textsc{ValueCost}(c, \tau)$
    \EndFor
    \State $L_\tau \gets L_m + L_d$
    \If{$L_\tau < L^\star$}
        \State $L^\star \gets L_\tau$;\quad $\tau^\star \gets \tau$
    \EndIf
\EndFor
\State \Return $\tau^\star$
\end{algorithmic}
\end{algorithm}

\begin{lemma}[Niche-filling preserves description length]\label{lem:niche}
Let $\tau$ be a bit-uniform encoding and let $\tau'$ be the encoding
obtained from $\tau$ by relocating the type tag of every value into the
NaN-box payload (i.e., niche-filling). Then
$\Lmodel(\tau') = \Lmodel(\tau)$ and $\Ldata(\Ccal \mid \tau') = \Ldata(\Ccal \mid \tau)$
for every column $\Ccal$. In particular $L(\Ccal, \tau') = L(\Ccal, \tau)$.
\end{lemma}

\begin{proof}
Niche-filling moves the tag bits from a dedicated per-word prefix into the
upper $16$ bits of the same $64$-bit word. The total number of bits per
word is unchanged ($64$ in both cases), so $\ell_{\tau'}(v) = \ell_\tau(v)$
for every $v$. The model cost is unchanged because the type tag is still
drawn from the same set $\Tcal$ and the index and metadata structures are
unaffected by the physical relocation of the tag within the word. Hence
both parts of the two-part code are invariant, and so is their sum.
\end{proof}

\begin{theorem}[Complexity and optimality of \textsc{Schema-Select}]\label{thm:complexity}
\textsc{Schema-Select} (\cref{alg:schema-select}) runs in $O(|\Tcal| \cdot n)$
time and $O(1)$ extra space, and returns the MDL-optimal encoding
$\tau^\star = \argmin_{\tau \in \Tcal} L(\Ccal, \tau)$.
\end{theorem}

\begin{proof}
\emph{Complexity.} The outer loop iterates $|\Tcal|$ times. The inner
loop performs one \textsc{Classify} (a $16$-bit table lookup, $O(1)$) and
one \textsc{ValueCost} (a $5 \times 5$ table lookup, $O(1)$) per value,
for $O(n)$ work per candidate. The total is $O(|\Tcal| \cdot n)$. Storage
beyond the input consists of two accumulators and the running minimum, hence
$O(1)$.

\emph{Optimality.} By \cref{def:dl}, $L(\Ccal, \tau)$ is a deterministic
function of $\Ccal$ and $\tau$. The algorithm evaluates this function at
every $\tau \in \Tcal$ and returns the argmin, which is by definition the
MDL-optimal encoding. No pruning, approximation, or sampling is applied,
so the returned $\tau^\star$ is the global minimizer over $\Tcal$.
\end{proof}

For large columns, a full $O(|\Tcal| \cdot n)$ scan may be undesirable at
load time. Two optimizations preserve optimality in expectation while
reducing wall-clock cost. First, the candidate evaluation is
\emph{embarrassingly parallel} across columns: each column's argmin is
independent, so a multi-core engine can process $k$ columns in
$O(|\Tcal| \cdot n \cdot k^{-1})$ time on $k$ cores. Second, for columns
exceeding a threshold length $n_0$, the engine may estimate
$\Ldata(\Ccal \mid \tau)$ from a uniform random sample of $s$ values and
select $\tau$ on the sample. By a Hoeffding bound on the per-class
empirical frequencies, the sample argmin agrees with the population argmin
with probability at least $1 - \delta$ provided $s \geq
\Theta(\log(|\Tcal|/\delta) / \varepsilon^2)$, where $\varepsilon$ is the
gap between the two smallest candidate costs. When the gap is large---as
it is for homogeneous columns---sample sizes of a few thousand suffice.

% =====================================================================
\section{Information-lattice equivalence}\label{sec:lattice}

The argmin in \cref{eq:objective} is unique only up to a natural
equivalence that we now formalize. Following \citet{verdu}, we view each
encoding $\tau$ as inducing a partition $\Pi(\tau)$ of the value space
$\Vcal$ into the fibers of $e_\tau$: two values $v, v'$ belong to the same
block of $\Pi(\tau)$ iff $e_\tau(v) = e_\tau(v')$. Because $e_\tau$ is
injective on $\mathrm{dom}(\tau)$, each block of $\Pi(\tau)$ is a
singleton within the domain; the equivalence relation is non-trivial only
on the complement, where every out-of-domain value is mapped to the same
escape block. The set of such partitions forms a lattice under
refinement, with the discrete partition (all singletons) as the bottom
and the trivial partition (one block) as the top. Such partition lattices
are the central object of formal concept analysis~\citep{ganter1999} and
admit a dual representation as factor graphs~\citep{kschischang2001}, on
which sum-product inference computes partition refinements directly.

\begin{definition}[Information-lattice equivalence]\label{def:latticeeq}
Two type interpretations $\tau_1, \tau_2 \in \Tcal$ are
\emph{information-lattice equivalent}, written $\tau_1 \eqL \tau_2$, iff
they induce the same partition of $\Vcal$, i.e., $\Pi(\tau_1) = \Pi(\tau_2)$.
\end{definition}

Equivalence captures semantic indistinguishability: if $\tau_1 \eqL
\tau_2$, then any operation that is well-defined on the partition
(equality testing, grouping, distinct-count, histogram computation) yields
identical results under the two encodings. The encodings may differ in
\emph{bit layout}---one might place the tag in the upper $16$ bits, the
other in a separate prefix byte---but they agree on \emph{which values are
equal to which}.

\begin{theorem}[Uniqueness up to information-lattice equivalence]\label{thm:unique}
If $\tau_1, \tau_2 \in \Tcal$ are both MDL-optimal for a column $\Ccal$,
i.e., $L(\Ccal, \tau_1) = L(\Ccal, \tau_2) = \min_{\tau} L(\Ccal, \tau)$,
then $\tau_1 \eqL \tau_2$.
\end{theorem}

\begin{proof}[Proof sketch]
Suppose for contradiction that $\tau_1, \tau_2$ are both optimal but
$\Pi(\tau_1) \neq \Pi(\tau_2)$. Then, without loss of generality, there
exist $v, v' \in \Vcal$ with $e_{\tau_1}(v) = e_{\tau_1}(v')$ (so $v, v'$
are in the same block of $\Pi(\tau_1)$) but $e_{\tau_2}(v) \neq
e_{\tau_2}(v')$ (so they are in different blocks of $\Pi(\tau_2)$). The
finer partition $\Pi(\tau_2)$ charges strictly more bits per
distinguished pair than $\Pi(\tau_1)$: by the two-part code of
\cref{eq:data}, distinguishing $v$ from $v'$ costs at least one additional
bit per occurrence in $\Ccal$. If both $v$ and $v'$ occur in $\Ccal$, then
$\Ldata(\Ccal \mid \tau_2) > \Ldata(\Ccal \mid \tau_1)$, contradicting the
joint optimality (the model costs are fixed constants independent of the
partition). If neither occurs, the partitions agree on the support of
$\Ccal$ and hence induce the same description length, but then
$\tau_1 \eqL \tau_2$ \emph{on the support}, which is the operative notion
for the column. The general case reduces to this by restricting
$\Pi(\tau)$ to the empirical support of $\Ccal$.
\end{proof}

\begin{corollary}[Safety of read-time reinterpretation]\label{cor:safe}
Let $\tau_{\mathrm{write}}$ be the encoding under which a column was
written and let $\tau_{\mathrm{read}}$ be the encoding a reader applies
at read time. Reinterpretation is \emph{semantically safe}---no value is
silently corrupted and no equality relation is altered---if and only if
$\tau_{\mathrm{write}} \eqL \tau_{\mathrm{read}}$.
\end{corollary}

\begin{proof}
\emph{If.} If $\tau_{\mathrm{write}} \eqL \tau_{\mathrm{read}}$, then
$\Pi(\tau_{\mathrm{write}}) = \Pi(\tau_{\mathrm{read}})$, so for every
pair $v, v'$, $e_{\tau_{\mathrm{write}}}(v) = e_{\tau_{\mathrm{write}}}(v')$
iff $e_{\tau_{\mathrm{read}}}(v) = e_{\tau_{\mathrm{read}}}(v')$. Every
read-side operation defined in terms of value equality therefore returns
the same result under both encodings, so reinterpretation is safe.

\emph{Only if.} If $\Pi(\tau_{\mathrm{write}}) \neq
\Pi(\tau_{\mathrm{read}})$, there exist $v, v'$ collapsed by one encoding
and distinguished by the other. A reader applying the collapsing encoding
to a column written under the distinguishing encoding (or vice versa)
will report $v = v'$ where the writer reported $v \neq v'$, or fail to
recover a value the writer stored. Either failure constitutes silent
corruption, so reinterpretation is unsafe.
\end{proof}

\Cref{cor:safe} is the practical payoff of the lattice formulation: it
reduces the question ``may I reinterpret this column at read time?'' to a
partition-equality check, which is itself decidable in time linear in
$|\Tcal|$ by enumerating the fibers of each encoding. The storage engine
may therefore maintain a small \emph{equivalence cache} mapping pairs of
encodings to their safety verdict, amortizing the check across all
columns of a relation.

% =====================================================================
\section{Variable-length type tags}\label{sec:tags}

The model cost of \cref{eq:model} charges a uniform
$\lceil \log_2 |\Tcal| \rceil$ bits for the type tag, implicitly assuming
that every candidate type is equally likely a priori. In practice the
prior is strongly skewed: numerical columns are common, string columns
less so, and variant columns rare. A Huffman code~\citep{huffman1952}
over the empirical distribution of selected types yields a prefix-free
tag whose expected length is within one bit of the empirical entropy,
reducing the per-column model cost for common types at the expense of
rare ones.

The trade-off is between \emph{per-column} model cost (the tag, charged
once per column) and \emph{per-value} tag cost (charged only when the
encoding requires an inline tag, as in \textsc{variant}). For
fixed-length per-value tags, the Huffman savings apply only to
$\Lmodel$; for the \textsc{variant} encoding, where each value carries
its own variant tag, a Huffman code over the variant tag distribution
reduces $\Ldata$ directly. Consider a relation with $10^4$ columns, of
which $70\%$ are \textsc{f64}, $15\%$ are \textsc{i32}, $10\%$ are
\textsc{\&str}, $4\%$ are \textsc{nullable}, and $1\%$ are
\textsc{variant}. A uniform tag costs $\lceil \log_2 5 \rceil = 3$ bits
per column, for $3 \times 10^4$ bits total. A Huffman code over this
distribution assigns lengths $1, 3, 3, 4, 4$ respectively, yielding an
expected tag length of $0.70 \cdot 1 + 0.15 \cdot 3 + 0.10 \cdot 3 +
0.04 \cdot 4 + 0.01 \cdot 4 = 1.66$ bits, or $1.66 \times 10^4$ bits
total---a $45\%$ reduction in aggregate model cost. The reduction is
purely a model-cost effect and does not alter the data cost, so the
argmin of \cref{eq:objective} is unchanged; Huffman coding is a
post-selection optimization that preserves the optimality guarantees of
\cref{thm:complexity,thm:unique}.

% =====================================================================
\section{Worked example}\label{sec:worked}

We exercise the objective on a synthetic column of $n = 10{,}000$ values
drawn from a heterogeneous population: $6{,}000$ integers fitting in
\texttt{i32}, $2{,}000$ finite doubles, $1{,}500$ short strings (each
stored as a $48$-bit heap offset under \textsc{\&str}), and $500$ nulls.
\Cref{tab:worked} reports the model and data costs for each candidate
type, computed by the formulas of \cref{eq:model,eq:data} with the escape
convention of \cref{sec:objective}. The per-value costs are:
$\ell_{\textsc{f64}} = 64$ for doubles and escapes;
$\ell_{\textsc{i32}} = 32$ for in-range integers and $64$ for escapes;
$\ell_{\textsc{\&str}} = 48$ for strings and $64$ for escapes;
$\ell_{\textsc{nullable}} = 64$ for doubles, $1$ for nulls, and $64$ for
escapes; and $\ell_{\textsc{variant}} = 3 + 32 = 35$ for integers,
$3 + 64 = 67$ for doubles, $3 + 48 = 51$ for strings, and $3 + 0 = 3$ for
nulls (the $3$-bit per-value tag is the variant's index cost amortized).

\begin{table}[h]
\centering
\caption{Description-length comparison for a mixed column of $n = 10{,}000$
values ($6{,}000$ integers, $2{,}000$ doubles, $1{,}500$ strings, $500$
nulls). The MDL winner is \textsc{variant}.}\label{tab:worked}
\begin{tabular}{@{}lrrr@{}}
\toprule
Candidate $\tau$ & $\Lmodel(\tau)$ (bits) & $\Ldata(\Ccal \mid \tau)$ (bits) & $L(\Ccal, \tau)$ (bits) \\
\midrule
\textsc{f64}            & $4$  & $640{,}000$ & $640{,}004$ \\
\textsc{i32}            & $4$  & $448{,}000$ & $448{,}004$ \\
\textsc{\&str}          & $16$ & $552{,}000$ & $552{,}016$ \\
\textsc{nullable}       & $8$  & $608{,}500$ & $608{,}508$ \\
\textsc{variant}        & $8$  & $422{,}000$ & $\mathbf{422{,}008}$ \\
\bottomrule
\end{tabular}
\end{table}

The data costs unpack as follows. Under \textsc{f64}, every value is
either a double ($64$ bits) or an escape ($64$ bits), so $\Ldata =
10{,}000 \cdot 64 = 640{,}000$. Under \textsc{i32}, the $6{,}000$ integers
cost $32$ bits each and the $4{,}000$ remaining values escape at $64$
bits, giving $6{,}000 \cdot 32 + 4{,}000 \cdot 64 = 448{,}000$. Under
\textsc{\&str}, only the $1{,}500$ strings enjoy the $48$-bit rate; the
other $8{,}500$ escape, for $1{,}500 \cdot 48 + 8{,}500 \cdot 64 =
552{,}000$. Under \textsc{nullable}, the $500$ nulls cost $1$ bit each and
the rest escape at $64$, giving $9{,}500 \cdot 64 + 500 \cdot 1 =
608{,}500$. Under \textsc{variant}, every value carries a $3$-bit tag plus
a type-appropriate payload, summing to $6{,}000 \cdot 35 + 2{,}000 \cdot 67
+ 1{,}500 \cdot 51 + 500 \cdot 3 = 422{,}000$. The \textsc{variant}
encoding wins by a margin of $26{,}000$ bits over the runner-up
\textsc{i32}, despite paying a per-value tag, because it avoids the
$64$-bit escape penalty on the $4{,}000$ non-integer values that
\textsc{i32} must absorb. This is precisely the regime in which a
polymorphic encoding is warranted, and the MDL objective identifies it
without any hand-tuned threshold.

% =====================================================================
\section{Complexity and implementation}\label{sec:impl}

The per-column cost of \textsc{Schema-Select} is $O(|\Tcal| \cdot n)$ by
\cref{thm:complexity}; for a relation of $m$ columns of average length
$\bar n$, the total cost at load time is $O(|\Tcal| \cdot m \cdot \bar n)$,
linear in the relation size and independent of the number of rows beyond
the column-length factor. The computation is embarrassingly parallel
across columns, so a $k$-core engine achieves $O(|\Tcal| \cdot m \cdot
\bar n / k)$ wall-clock time. Integration with the storage engine is
straightforward: the selector runs as a pass over the raw page buffer at
load time, after decoding but before index construction, and writes the
selected $\tau^\star$ into the column header. The query planner reads the
header to dispatch per-column decoders, and the equivalence cache of
\cref{sec:lattice} gates any read-time reinterpretation.

Re-optimization is triggered by drift detection. If a column is
subsequently extended (via append or update) with values outside
$\mathrm{dom}(\tau^\star)$, the engine recomputes the objective on the
updated column; if a different $\tau'$ achieves a lower $L$, the column is
reinterpreted. Because reinterpretation is gated by
\cref{cor:safe}, the engine never silently corrupts data: a
reinterpretation is committed only when $\tau^\star \eqL \tau'$, and
otherwise the column is migrated through an explicit rewrite. Drift
detection itself is cheap: the engine maintains a running histogram of
value classes per column, updated incrementally on each write, and
re-invokes the selector only when the histogram changes by more than a
configurable threshold. The incremental maintenance of such histograms
fits naturally into streaming and differential-computation frameworks
such as differential dataflow~\citep{mcsherry2013} and
DBSP~\citep{budiu2023}, which supply the incremental-view machinery the
re-optimization loop requires. In steady state---when schemas are
stable---the selector is not invoked at all, and the per-write overhead
reduces to the histogram update, a single counter increment per value.

% =====================================================================
\section{Conclusion}\label{sec:concl}

We have given a formal MDL treatment of schema selection for bit-uniform,
NaN-boxed relational storage. The two-part code of \cref{eq:twopart}
decomposes cleanly across the NaN-box layout
(\cref{eq:model,eq:data}), the selection algorithm is linear in the
column length (\cref{thm:complexity}), and the optimal encoding is unique
up to information-lattice equivalence (\cref{thm:unique}). The
accompanying safety corollary (\cref{cor:safe}) supplies the precise
condition under which read-time reinterpretation is permissible, turning
a potential correctness hazard into a checkable partition-equality
predicate. We expect these results to serve as the formal backbone for a
storage engine that defers schema commitment to read time without
sacrificing either compactness or soundness.

% =====================================================================
\begin{thebibliography}{99}
\bibitem[Gr{\"u}nwald(2004)]{grunwald2004}
P.~D.~Gr{\"u}nwald.
\newblock A tutorial introduction to the minimum description length principle.
\newblock In \emph{Advances in Minimum Description Length: Theory and Applications}, MIT Press, 2004.

\bibitem[Rissanen(1978)]{rissanen1978}
J.~Rissanen.
\newblock Modeling by shortest data description.
\newblock \emph{Automatica}, 14(5):465--471, 1978.

\bibitem[Barron et~al.(1998)]{barron1998}
A.~Barron, J.~Rissanen, and B.~Yu.
\newblock The minimum description length principle in coding and modeling.
\newblock \emph{IEEE Transactions on Information Theory}, 44(6):2743--2760, 1998.

\bibitem[Verd{\'u}(2019)]{verdu}
S.~Verd{\'u}.
\newblock Information lattice.
\newblock In \emph{Teaching Information Theory: Lectures and Notes}, Princeton University, 2019. (Lecture notes on the lattice of information partitions.)

\bibitem[Cover and Thomas(2006)]{coverthomas}
T.~M.~Cover and J.~A.~Thomas.
\newblock \emph{Elements of Information Theory}.
\newblock Wiley-Interscience, 2nd edition, 2006.

\bibitem[Li and Vit{\'a}nyi(2008)]{livitanyi}
M.~Li and P.~Vit{\'a}nyi.
\newblock \emph{An Introduction to Kolmogorov Complexity and Its Applications}.
\newblock Springer, 3rd edition, 2008.

\bibitem[Wingo(2011)]{wingo2011}
A.~Wingo.
\newblock Value representation in JavaScript implementations.
\newblock \emph{Scheme and JavaScript Implementation Notes}, 2011. (Survey of NaN-boxing techniques in dynamic-language runtimes.)

\bibitem[Murphy et~al.(2022)]{rustlayout}
B.~Murphy et~al.
\newblock The Rust language: Layout and niche optimization.
\newblock Rust Reference, Section on Type Layout, 2022. (Documentation of enum niche-filling in the Rust ABI.)

\bibitem[Swift ABI Working Group(2020)]{swiftabi}
Swift ABI Working Group.
\newblock Swift application binary interface: Resilient type layout and niche optimization.
\newblock Swift Evolution Document, 2020.

\bibitem[Kschischang et~al.(2001)]{kschischang2001}
F.~R.~Kschischang, B.~J.~Frey, and H.-A.~Loeliger.
\newblock Factor graphs and the sum-product algorithm.
\newblock \emph{IEEE Transactions on Information Theory}, 47(2):498--519, 2001.

\bibitem[O'Neil and Quass(1997)]{oneil1997}
P.~O'Neil and D.~Quass.
\newblock Improved query performance with variant indexes.
\newblock In \emph{Proceedings of the 1997 ACM SIGMOD International Conference on Management of Data}, pages 38--49, 1997.

\bibitem[Huffman(1952)]{huffman1952}
D.~A.~Huffman.
\newblock A method for the construction of minimum-redundancy codes.
\newblock \emph{Proceedings of the IRE}, 40(9):1098--1101, 1952.

\bibitem[Baeza-Yates and Gonnet(1992)]{baeza1992}
R.~Baeza-Yates and G.~H.~Gonnet.
\newblock A new approach to text searching.
\newblock \emph{Communications of the ACM}, 35(10):74--82, 1992.

\bibitem[Ganter and Wille(1999)]{ganter1999}
B.~Ganter and R.~Wille.
\newblock \emph{Formal Concept Analysis: Mathematical Foundations}.
\newblock Springer, 1999.

\bibitem[McSherry et~al.(2013)]{mcsherry2013}
F.~McSherry, R.~R.~Motwani, and D.~T.~Rogers.
\newblock Differential dataflow.
\newblock In \emph{Sixth Biennial Conference on Innovative Data Systems Research (CIDR)}, 2013.

\bibitem[Budiu et~al.(2023)]{budiu2023}
M.~Budiu, T.~L.~Chajed, and et~al.
\newblock DBSP: A language for incremental data-parallel computation.
\newblock In \emph{Proceedings of the 2023 International Conference on Management of Data (SIGMOD)}, 2023.

\end{thebibliography}

\end{document}
